\documentclass[11pt]{article}
\usepackage[english,russian]{babel}
\usepackage[utf8]{inputenc}
\usepackage[T2A]{fontenc}
\usepackage{url}
\usepackage{graphicx,DCCN2026_ru}

\pagestyle{fancy}
\fancyhead{}
\fancyfoot{}

\linespread{1.0}

\usepackage{amsmath,amssymb}
\usepackage{booktabs}
\usepackage{tikz}
\usetikzlibrary{positioning,arrows.meta,fit,backgrounds,calc}

\makeatletter
\fancyhead[RO]{\small DCCN 2026\\ {21--25 сентября 2026}}
\fancyhead[LO]{\small Без авторов (double-blind)\\
Иерархическая когнитивная архитектура управления}

\c@page=1
\makeatother


\title{Иерархическая когнитивная архитектура управления для мобильной робототехнической платформы}

\author[1]{\small }
\author[2]{\small }

\email{ }


% Compact spacing
\setlength{\textfloatsep}{4pt plus 1pt minus 1pt}
\setlength{\floatsep}{4pt plus 1pt minus 1pt}
\setlength{\intextsep}{4pt plus 1pt minus 1pt}
\setlength{\abovedisplayskip}{3pt plus 1pt minus 1pt}
\setlength{\belowdisplayskip}{3pt plus 1pt minus 1pt}
\setlength{\abovedisplayshortskip}{0pt plus 1pt}
\setlength{\belowdisplayshortskip}{2pt plus 1pt minus 1pt}


% Even tighter spacing for 6p version
\setlength{\parskip}{0pt}
\setlength{\parindent}{1em}
\begin{document}

\udc{004.896}

{\let\newpage\relax\maketitle}
\vskip -1.5em
%\footnotetext{Работа выполнена при поддержке ...}


\begin{abstract}
Предложена реализация трёхуровневой иерархической когнитивной архитектуры управления мобильным агентом с дифференциальным приводом в среде с препятствиями и переменными свойствами поверхности, объединяющая нижний уровень (TD3), средний (LQR) и верхний (обход по карте). Покомпонентный анализ показывает, что верхний уровень разгружает нижний при известной геометрии препятствий, но не компенсирует физических ограничений нижнего уровня при слабом сцеплении колёс.

\keywords{обучение с подкреплением, TD3, LQR, иерархическое управление,
мобильная робототехника, STRL, residual policy, когнитивный агент}
\end{abstract}


\section{Введение}
\label{sec:intro}

Автономная навигация мобильного агента в среде с препятствиями и переменными свойствами поверхности объединяет обучение с подкреплением, классическое управление и иерархическое планирование. Один из устоявшихся подходов~--- трёхуровневая иерархическая когнитивная архитектура~\cite{gat1998three,kiselev2020dissertation,kiselev2020} с распределением поведения по горизонтам планирования. В диссертации~\cite{kiselev2020dissertation} такая архитектура реализована для платформы МП-РМ <<Зарница>> с $Q$-обучением на нижнем уровне; в~\cite{veitsenfeld2024} направление развивается мультимодальной модельной компонентой.

В настоящей работе предложена реализация трёхуровневой архитектуры для \emph{непрерывного} управления: нижний уровень~--- стратегия, обученная методом TD3~\cite{fujimoto2018td3}; средний~--- LQR-регулятор; верхний~--- модуль обхода по известной карте. Архитектура протестирована в PyBullet на платформе Husky.

\textbf{Вклад работы.} (1)~Реализация трёхуровневой иерархической архитектуры для непрерывного управления агентом. (2)~Покомпонентный анализ распределения функций между уровнями на 65 конфигурациях ($7 \times 5 + 6 \times 5$) с $n = 10$ запусками на ячейку. (3)~Открытый репозиторий с воспроизводимой реализацией.


\section{Обзор работ}
\label{sec:related}

В задачах непрерывного управления применяются методы актора-критика~\cite{lillicrap2016ddpg,fujimoto2018td3,haarnoja2018sac,schulman2017ppo}; в настоящей работе используется TD3~\cite{fujimoto2018td3} с эталонной реализацией в \verb|stable-baselines3|~\cite{raffin2021sb3}. Комбинирование обученной стратегии с классическим регулятором~\cite{silver2018residual,johannink2019residual} компенсирует ограничения каждого подхода в отдельности. Параллельное направление гибридизации обучения с подкреплением и методов оптимального управления систематизировано в обзорах~\cite{hewing2020lbmpc,brunke2022safe}, с близкими подходами в~\cite{levine2016end,cheng2019control,kahn2018selfsup}. Настоящая работа развивает направление STRL~\cite{gat1998three,kiselev2020dissertation,kiselev2020,veitsenfeld2024} в части непрерывного управления и количественного анализа роли каждого уровня архитектуры.

\section{Метод}
\label{sec:method}

\subsection{Постановка задачи}
\label{sec:problem}

Задача формализуется как марковский процесс принятия решений $\langle S, A, T, R, \gamma \rangle$ с непрерывными пространствами в плоской рабочей области $10 \times 10$ м с препятствиями известной геометрии. Среда моделируется в PyBullet~\cite{coumans2021pybullet} на платформе Husky. Состояние агента
\begin{equation}
s_t = (\mathbf{r}_t,\ \theta_t,\ \mathbf{v}_t,\ \omega_t,\ \mathbf{g}_t,\ \ell_t)^\top \in \mathbb{R}^{n_s},
\end{equation}
включает координаты $\mathbf{r}_t = (x_t, y_t)$, ориентацию $\theta_t$, линейную и угловую скорости $\mathbf{v}_t, \omega_t$, относительное положение цели $\mathbf{g}_t = \mathbf{r}_g - \mathbf{r}_t$, и лидарные измерения $\ell_t \in \mathbb{R}^{16}$ (только при наличии препятствий). Действие $a_t \in [-1, 1]^2$~--- нормированные скорости левой и правой пар колёс, преобразуемые в физические $\omega_w \in [-10, 10]$ рад/с. Шаги физики и управления $1/240$ и $1/20$ с, $\gamma = 0.99$. Цель на расстоянии 2--5 м, радиус достижения 0.4 м. Эпизод завершается при достижении цели, перевороте, выходе за пределы рабочей области, столкновении или таймауте 500 шагов. Вознаграждение плотное (штраф за расстояние, шаговый штраф, терминальные бонус и штрафы).


\subsection{Иерархическая архитектура управления}
\label{sec:architecture}



Архитектура , восходящая к трёхуровневой схеме~\cite{gat1998three} и развитая для когнитивных агентов в~\cite{kiselev2020dissertation,kiselev2020}адаптирована для непрерывного управления и включает три уровня.

\textit{Нижний уровень}~--- обученная методом TD3 стратегия $\pi_{\mathrm{TD3}}(s) \to a$, опирающаяся на текущее наблюдение.

\textit{Средний уровень}~--- LQR-регулятор по линеаризованной модели поперечного отклонения и отклонения курса агента от прямой к цели. Регулятор выдаёт скалярную коррекцию $u_{\mathrm{LQR}} \in [-u_{\max}, u_{\max}]$ к угловой скорости, комбинируемую с действием нижнего уровня по схеме остаточной коррекции~\cite{silver2018residual,johannink2019residual}.

\textit{Верхний уровень}~--- модуль обхода препятствий по известной карте (MapAware) с индикатором активации $\mathbb{1}_{\mathrm{map}}(s) \in \{0, 1\}$. При приближении агента к препятствию в передней полусфере на расстояние ниже $d_{\mathrm{takeover}} = 1.5$ м индикатор переключается в 1, при выходе за $d_{\mathrm{release}} = 2.5$ м~--- в 0; гистерезис предотвращает дребезг. Пороги согласованы с кинематикой: $v_{\max} = \omega_{\max} \cdot r_w \approx 1.8$ м/с при $\omega_{\max} = 10$ рад/с и $r_w = 0.178$ м, тогда $d_{\mathrm{takeover}} / v_{\max} \approx 0.83$ с обеспечивает запас для смены курса. В текущей реализации уровень замещает управление в опасных конфигурациях, не предлагая последовательности подцелей; полноценный планировщик~--- естественное направление развития.

Конечное действие формируется как $a_t = a_{\mathrm{MapAware}}(s_t)$ при $\mathbb{1}_{\mathrm{map}}(s_t) = 1$, иначе как $\mathrm{clip}(a_{\mathrm{TD3}}(s_t) + \alpha [-1,+1]^\top u_{\mathrm{LQR}}(s_t), -1, 1)$, где $\mathbb{1}_{\mathrm{map}}$~--- индикатор активности верхнего уровня, $\alpha = 0.2$~--- весовой коэффициент остаточной коррекции.


\subsection{Нижний уровень: TD3}
\label{sec:reactive}

В качестве базовой стратегии используется TD3~\cite{fujimoto2018td3}, актор-критик с двойной оценкой функции ценности, задержанным обновлением актора и буфером переходов; функции ценности и стратегия~--- полносвязные сети со слоями 400 и 300 (ReLU). Гиперпараметры стандартные для TD3~\cite{raffin2021sb3}: минибатч 256, Adam с lr $10^{-3}$, $\tau = 0.005$, $d = 2$, $\sigma_t = 0.2$, $c = 0.5$, шум $\mathcal{N}(0, 0.3)$, буфер $|\mathcal{D}| = 2 \cdot 10^5$.

Вектор наблюдения $o_t = (\cos\theta_t, \sin\theta_t, \mathbf{v}_t, \omega_t, \mathbf{g}_t, \ell_t)^\top$ включает ориентацию, скорости и относительное положение цели $\mathbf{g}_t$. Лидарные измерения $\ell_t \in \mathbb{R}^{16}$ получаются агрегированием 360 лучей (паспортные характеристики LDS-01 TurtleBot~3 Burger~\cite{lds01_datasheet}) в 16 секторов по $22.5^\circ$ операцией поиска минимума $\ell_t^{(k)} = \min_{j \in \mathcal{B}_k} d_t^{(j)}$. Такая агрегация гарантирует обнаружение тонких препятствий, которые могут проваливаться между лучами в низкоразрешённом сканировании. Наблюдение привилегированное, что отделяет управление от восприятия.


\subsection{Средний уровень: LQR-residual}
\label{sec:tactical}

LQR-регулятор работает с вектором ошибки $e = [e_y, e_\theta]^\top$ (поперечное отклонение и отклонение курса от прямой к цели) и выдаёт коррекцию $\omega = -K e$ с матрицей $K$, найденной из решения алгебраического уравнения Риккати для линеаризованной модели c $Q = \mathrm{diag}(2.0, 1.0)$ и $R = 5.0$ (преобладание штрафа на управление обеспечивает мягкие коррекции).

В реализации $e_y$ обнуляется на каждом шаге: опорная прямая переопределяется от текущей позы агента к цели, и закон вырождается в $\omega = -k_\theta e_\theta$, P-регулятор по углу с коэффициентом из LQR-задачи. Выход ограничивается $|\omega| \leq u_{\max} = 0.3$.


\subsection{Верхний уровень: обход с использованием карты}
\label{sec:supervisory}

Уровень получает карту препятствий в виде списка $\{(o^x_i, o^y_i, r_i)\}$ и вычисляет расстояния до краёв препятствий в передней полусфере $\pm \pi/2$ относительно курса агента. Активация определяется двупороговой логикой с гистерезисом: индикатор $\mathbb{1}_{\text{map}}$ переключается в 1 при пересечении $d_{\text{takeover}}$ и обратно в 0 при пересечении $d_{\text{release}}$. При активном уровне управление задаётся прямолинейным движением со скоростью $v_{\mathrm{sup}} = 0.4$ м/с и угловой скоростью $\omega_{\mathrm{sup}} = k_{\mathrm{turn}} \cdot \mathrm{sign}(-\theta_{\mathrm{nearest}})$, $k_{\mathrm{turn}} = 1.5$, где $\theta_{\mathrm{nearest}}$~--- угол на ближайшее препятствие. Простота алгоритма преднамеренна: она позволяет изолировать модельное знание о среде. Для сопоставления модельного и сенсорного знания используется также альтернативный модуль обхода по лидарным дальностям (LidarAvoid).


\section{Эксперименты}
\label{sec:experiments}

\subsection{Конфигурация и метрики}
\label{sec:exp_setup}

Эксперименты проведены на сервере 8 ядер Ryzen 9 9950X / 16 ГБ RAM без GPU, PyBullet 3.2 в безголовом режиме, 6 параллельных сред через \verb|SubprocVecEnv|, сиды $\{200, \ldots, 209\}$, $n = 10$ запусков на ячейку. Основная метрика~--- доля успешных эпизодов (success rate). Дополнительно фиксируются длина эпизода, гладкость управления, максимальная линейная скорость и доля времени в верхнеуровневом режиме. Сравниваются пять архитектурных режимов: TD3, TD3+LQR, TD3+LidarAvoid, TD3+MapAvoid и TD3+MapAvoid+LQR.


\subsection{Эксперимент A: множественные конфигурации препятствий}
\label{sec:exp_a}

Семь детерминированных конфигураций препятствий: одиночное препятствие на прямой, два со смещением, коридор из трёх, барьер с проходом, слалом, диагональное расположение цели, широкое препятствие.

\begin{table}[h]
\begin{center}
\begin{tabular}{lccccc}
\toprule
Сценарий & TD3 & +LQR & +Lidar & +Map & +Map+LQR \\
\midrule
single\_on\_line     & 0/10 & 0/10 & \textbf{10/10} & \textbf{10/10} & \textbf{10/10} \\
two\_offset         & 0/10 & 0/10 & \textbf{10/10} & \textbf{10/10} & \textbf{10/10} \\
three\_corridor     & 0/10 & 0/10 & 0/10           & \textbf{10/10} & \textbf{10/10} \\
barrier\_with\_gap  & \textbf{10/10} & \textbf{10/10} & \textbf{10/10} & 0/10 & 0/10 \\
slalom              & 0/10 & 0/10 & 0/10 & 0/10           & 0/10           \\
diagonal\_goal      & 10/10 & 10/10 & 10/10 & 10/10        & 10/10         \\
wide\_obstacle      & 0/10 & 0/10 & 0/10 & \textbf{10/10} & \textbf{10/10} \\
\bottomrule
\end{tabular}
\caption{Эксперимент~A: доля успешных эпизодов (из 10 запусков на сидах
200--209) для семи конфигураций препятствий и пяти архитектурных режимов.
Все ячейки 100\,\%~детерминированы.}
\label{tab:exp_a}
\end{center}
\end{table}

Сценарий \emph{diagonal\_goal} без препятствий на прямой решается всеми режимами и служит контролем. Сценарий \emph{slalom} не решается ни одним режимом, фиксируя границу применимости архитектуры.

В сценариях \emph{single\_on\_line} и \emph{two\_offset} оба варианта верхнего уровня обеспечивают полный успех, тогда как нижний и средний уровни в одиночку не справляются: обученная стратегия идёт по прямой к цели, а LQR-коррекция общую стратегию не меняет. В \emph{three\_corridor} полный успех достигает только модельный верхний уровень, что соответствует роли стратегического уровня при известной структуре среды. Обратный паттерн наблюдается в \emph{barrier\_with\_gap}, где модельный уровень обходит барьер целиком, хотя геометрический проход проходим напрямую, и в \emph{wide\_obstacle}, где аппроксимация цилиндром малого радиуса не активирует MapAvoid на нужной дистанции. Эти сценарии указывают, что полезность каждого варианта верхнего уровня зависит от соответствия внутренней модели реальной геометрии среды.


\subsection{Эксперимент B: переменные свойства поверхности}
\label{sec:exp_b}

Реализованы три типа поверхности: \emph{NORMAL} (трение 0.9), \emph{ICE} (0.05) и \emph{SAND} (0.6 с дополнительным сопротивлением качению), каждый в конфигурациях \emph{empty} и \emph{with\_obstacle}, $n = 10$.

\begin{table}[h]
\begin{center}
\begin{tabular}{llccccc}
\toprule
Поверхность & Конфиг. & TD3 & +LQR & +Lidar & +Map & +Map+LQR \\
\midrule
NORMAL & empty     & 10/10 & 10/10 & 10/10 & 10/10 & 10/10 \\
NORMAL & + obst.   & 0/10  & 0/10  & 0/10  & \textbf{10/10} & \textbf{10/10} \\
ICE    & empty     & 10/10 & 10/10 & 10/10 & 10/10 & 10/10 \\
ICE    & + obst.   & 0/10  & 0/10  & 0/10  & 0/10  & \textbf{10/10} \\
SAND   & empty     & 10/10 & 10/10 & 10/10 & 10/10 & 10/10 \\
SAND   & + obst.   & 0/10  & 0/10  & \textbf{10/10}  & \textbf{10/10} & \textbf{10/10} \\
\bottomrule
\end{tabular}
\caption{Эксперимент~B: доля успешных эпизодов (из 10) для трёх типов поверхности и двух конфигураций.}
\label{tab:exp_b}
\end{center}
\end{table}

В отсутствие препятствий нижний уровень успешен на всех поверхностях, автоматически адаптируя скорость. Принципиальный результат~--- сценарий ICE+obstacle, где полного успеха достигает только полная трёхуровневая конфигурация TD3+MapAvoid+LQR: без среднего уровня низкое трение приводит к проскальзыванию при попытке обхода. Это эмпирически подтверждает взаимодополнение уровней~--- ни один в отдельности не достаточен при физических ограничениях нижнего уровня.

\section{Обсуждение}
\label{sec:discussion}

\textit{Средний уровень не меняет стратегических решений нижнего уровня}~--- включение LQR не меняет долю успеха ни в одной ячейке таблицы~\ref{tab:exp_a}, но подавляет высокочастотные колебания курса. \textit{Верхний уровень разгружает нижний при адекватной геометрической модели}: в \emph{single\_on\_line} и \emph{two\_offset} верхний уровень обеспечивает полный успех; обратный эффект в \emph{barrier\_with\_gap} и \emph{wide\_obstacle} указывает на чувствительность к точности представления геометрии. \textit{Модельное знание не компенсирует физических ограничений}: на ICE+obstacle полный успех достигается только полной трёхуровневой архитектурой. \textit{Ограничения}: простейший геометрический алгоритм без оптимизационного планирования, привилегированное наблюдение нижнего уровня, отсутствие реальной платформы.

\section{Заключение}
\label{sec:conclusion}

Реализована и эмпирически исследована трёхуровневая иерархическая архитектура управления мобильным агентом с дифференциальным приводом~\cite{kiselev2020dissertation,kiselev2020}: нижний уровень~--- стратегия TD3, средний~--- LQR, верхний~--- модуль обхода по карте. Покомпонентный анализ на семи конфигурациях препятствий и трёх типах поверхности показывает, что верхний уровень разгружает нижний при известной геометрии, но не компенсирует слабое сцепление колёс; средний улучшает гладкость управления. Направления развития: верхний уровень с оптимизационным планировщиком, интеграция с мультимодальной компонентой~\cite{veitsenfeld2024}, перенос на реальную платформу с ROS.


{\scriptsize\begin{thebibliography}{18}
\bibitem{gat1998three}
Gat E. On three-layer architectures
// Artificial Intelligence and Mobile Robots / D.~Kortenkamp, R.~Bonasso, R.~Murphy (eds.). AAAI Press, 1998. P.~195--210.
\bibitem{kiselev2020dissertation}
Киселёв Г.~А. Интеллектуальная система планирования поведения коалиции робототехнических агентов со STRL архитектурой:
дис. \dots канд. техн. наук. Москва, 2020. 117~с.
\bibitem{kiselev2020}
Киселёв Г.~А., Панов А.~И. Иерархическое психологически-инспирированное планирование для интеллектуальных динамических систем
// Известия РАН. Теория и системы управления. 2020. № 1. С.~123--137.
\bibitem{veitsenfeld2024}
Вейценфельд Д.~А., Киселёв Г.~А. Метод синтеза поведения когнитивного агента на основе обработки мультимодальных сигналов
// Моделирование и анализ данных. 2024. Т.~14. № 4. С.~45--62.
\bibitem{fujimoto2018td3}
Fujimoto S., van Hoof H., Meger D. Addressing function approximation error in actor-critic methods
// Proc. of ICML. 2018. P.~1587--1596.
\bibitem{lillicrap2016ddpg}
Lillicrap T., Hunt J., Pritzel A. et al. Continuous control with deep reinforcement learning
// Proc. of ICLR. 2016.
\bibitem{haarnoja2018sac}
Haarnoja T., Zhou A., Abbeel P., Levine S. Soft actor-critic: Off-policy maximum entropy deep reinforcement learning with a stochastic actor
// Proc. of ICML. 2018. P.~1861--1870.
\bibitem{schulman2017ppo}
Schulman J., Wolski F., Dhariwal P., Radford A., Klimov O. Proximal policy optimization algorithms
// arXiv preprint arXiv:1707.06347. 2017.
\bibitem{raffin2021sb3}
Raffin A., Hill A., Gleave A. et al. Stable-Baselines3: Reliable reinforcement learning implementations
// Journal of Machine Learning Research. 2021. V.~22. № 268. P.~1--8.
\bibitem{silver2018residual}
Silver T., Allen K., Tenenbaum J., Kaelbling L. Residual policy learning
// arXiv preprint arXiv:1812.06298. 2018.
\bibitem{johannink2019residual}
Johannink T., Bahl S., Nair A. et al. Residual reinforcement learning for robot control
// Proc. of ICRA. 2019. P.~6023--6029.
\bibitem{hewing2020lbmpc}
Hewing L., Wabersich K. P., Menner M., Zeilinger M. N. Learning-Based Model Predictive Control: Toward Safe Learning in Control // Annual Review of Control, Robotics, and Autonomous Systems. 2020. V.~3. P.~269--296.
\bibitem{brunke2022safe}
Brunke L., Greeff M., Hall A. W., Yuan Z., Zhou S., Panerati J., Schoellig A. P. Safe Learning in Robotics: From Learning-Based Control to Safe Reinforcement Learning // Annual Review of Control, Robotics, and Autonomous Systems. 2022. V.~5. P.~411--444.
\bibitem{levine2016end}
Levine S., Finn C., Darrell T., Abbeel P. End-to-end training of deep visuomotor policies
// Journal of Machine Learning Research. 2016. V.~17. P.~1--40.
\bibitem{cheng2019control}
Cheng R., Orosz G., Murray R. M., Burdick J. W. End-to-end safe reinforcement learning through barrier functions for safety-critical continuous control tasks
// Proc. of AAAI. 2019. P.~3387--3395.
\bibitem{kahn2018selfsup}
Kahn G., Villaflor A., Ding B., Abbeel P., Levine S. Self-supervised deep reinforcement learning with generalized computation graphs for robot navigation // Proc. of IEEE ICRA. 2018. P.~5129--5136.
\bibitem{coumans2021pybullet}
Coumans E., Bai Y. PyBullet, a Python module for physics simulation for games, robotics and machine learning. 2016--2021. URL: \url{http://pybullet.org}
\bibitem{lds01_datasheet}
Robotis Co., Ltd. LDS-01 (360 Laser Distance Sensor): Spec sheet // URL: \url{https://emanual.robotis.com/docs/en/platform/turtlebot3/appendix_lds_01/}. 2017.
\end{thebibliography}}

\end{document}